\section{The two model versions, and simplified waste schemes}\label{sec:sp:versions}

Sections~\ref{sec:sp:problem} and~\ref{sec:sp:solution} derived and read the planner's problem under
Version B, the closure in which the material intensity of new capital goods, $\Phi^I(t)$, is
exogenous. This section does two things. Section~\ref{sec:sp:versions:A} re-derives the problem under
Version A, $\Phi^I=\Omega\phi^I$, and shows that the entire difference between the two closures can be
written as a single scalar wedge applied at every point where the final good is used. Section
\ref{sec:sp:versions:simple} then catalogues the simplifications of the waste-generation scheme that
the rest of the note and the quantitative implementation make use of, and --- the point of the
exercise --- classifies them by whether they touch the allocation or only the reported mass flows.
The classification is different under the two versions, and that difference is the sharpest reason
to carry both.


\subsection{Version A: the material intensity of capital tied to the final good}
\label{sec:sp:versions:A}

\smalltitle{The closure.} Under Version A the material content of capital goods per unit of
expenditure is not pinned down independently of the material content of final-good activity in
general. Capital goods are made from the same composite good as everything else, so
$\Phi^I=\Omega\phi^I$ with $\phi^I$ exogenous: capital-goods production is $\phi^I$ times as
material-intensive as final-good activity on average, $\phi^I=1$ being the parsimonious case. The
level $\Omega$ is then determined by~\eqref{eq:sp:setup:ybalance_OmegaPhi},
\begin{align}
R &= \Omega\left(\mathcal D+\phi^IG\right),
&
\mathcal D &= \omega^YY+\omega^{N}C^N+\omega^{D}C^D+\omega^WC^W+\omega^CC,
\label{eq:sp:versions:Omega-B}
\end{align}
which by~\eqref{eq:sp:setup:Omega-phi} delivers $\Omega\ge0$ automatically. Problem
\ref{prob:sp:planner} is modified in three ways: $\Omega$ joins the control
set,~\eqref{eq:sp:versions:Omega-B} joins the constraints with multiplier $\lambda\zeta$, $\Phi^I$ is
replaced by $\Omega\phi^I$ in the ledger~\eqref{eq:sp:problem:ledger} and in the law of motion for
$M^K$, and the mass constraint~\eqref{eq:sp:problem:mass} is deleted along with its multiplier
$\eta$. The ledger is no longer linear in $W$; it is the quadratic
\eqref{eq:sp:setup:ledger-quadratic}, with a unique positive root.

\smalltitle{The wedge.} Differentiating the modified Lagrangian with respect to $\Omega$ and using
$\mathcal D+\phi^IG=R/\Omega$ gives the multiplier in closed form. Define the share of material
input that is embodied in new capital,
\begin{align}
\chi &\equiv \frac{\Phi^IG}{R}\;\in\;\left[0,1\right),
\qquad\text{then}\qquad
\zeta = \chi\left(p^W-p^M\right),
\qquad
\Delta \equiv \zeta\Omega = \Omega\chi\left(p^W-p^M\right).
\label{eq:sp:versions:zeta}
\end{align}
$\Delta$ is the object that does all the work. Its interpretation is direct: spending one more unit
of the final good on any activity $j$ raises $\mathcal D$ by $\omega^j$, which
by~\eqref{eq:sp:versions:Omega-B} \emph{dilutes} $\Omega$ for given $R$, which lowers $\Phi^I$, which
means less of the incoming mass is embodied in new capital and more of it is diffused as waste
today. Under Version A the composition of final-good spending is therefore a genuine externality,
and $\Delta\omega^j$ is its size. Note that $\Delta$ carries the sign of $p^W-p^M$ --- the same
bracket that determines whether $q\lessgtr1$ in~\eqref{eq:sp:solution:q} --- so the wedge and the
capital-price wedge flip together, exactly as the two instruments in
Section~\ref{sec:sp:solution:signs} do.

\smalltitle{The conditions.} With $\Delta$ in hand the entire system can be written by substitution.
A unit of the final good \emph{spent} on activity $j$ has effective social cost
$\left(1+\Delta\omega^j\right)$, and a unit \emph{produced} has effective social value
$\left(1-\Delta\omega^Y\right)$ --- the asymmetry is not an error, since producing output itself adds
$\omega^YY$ to the waste base. Writing
\begin{align}
B^{\mathrm B} &\equiv \left(1-\Delta\omega^Y\right)F_R-p^W+d^Wp^P+\zeta,
\label{eq:sp:versions:B-value}
\end{align}
for the Version-B counterpart of $B$, the first-order conditions are
\begin{align}
\text{\emph{consumption:}}
&&
u'\!\left(C\right) &= \lambda\left(1+\Delta\omega^C\right),
\label{eq:sp:versions:foc:C}
\\
\text{\emph{investment:}}
&&
q &= 1-\Phi^I\left(1-\chi\right)\left(p^W-p^M\right),
\label{eq:sp:versions:foc:I}
\\
\text{\emph{extraction:}}
&&
B^{\mathrm B} &= \left(1+\Delta\omega^N\right)C^N_N+p^S+\Omega^{N,S}p^W+d^Wp^P,
\label{eq:sp:versions:foc:N}
\\
\text{\emph{exploration:}}
&&
p^S+p^X &= \left(1+\Delta\omega^D\right)C^D_D+\Omega^{D,S}p^W,
\label{eq:sp:versions:foc:D}
\\
\text{\emph{treatment:}}
&&
\alpha B^{\mathrm B}+\left(1-d^W\right)p^P &= \left(1+\Delta\omega^W\right)\frac{dc^T}{d\varpi},
\label{eq:sp:versions:foc:varpi}
\\
\text{\emph{recycling capital:}}
&&
\left(1-\Delta\omega^Y\right)F_K &= a'\!\left(x\right)B^{\mathrm B},
\label{eq:sp:versions:foc:KR}
\\
\text{\emph{waste:}}
&&
p^W &= \left(1+\Delta\omega^W\right)c^W\!\left(\varpi\right)
+\left(1-\varpi+d^W\varpi\right)p^P-\varpi\alpha B^{\mathrm B},
\label{eq:sp:versions:foc:W}
\end{align}
and the costate equations become
\begin{align}
r &= \frac{\left(1-\Delta\omega^Y\right)F_K+\dot q}{q}-\delta,
&
\dot p^S &= r\,p^S+\left(1+\Delta\omega^N\right)C^N_S,
&
\dot p^X &= r\,p^X+\left(1+\Delta\omega^D\right)C^D_X,
\notag\\
\dot p^P &= \Big[r+\theta+\theta'P\Big]p^P-d^{\mathrm B},
&
\dot p^M &= \left(r+\delta\right)p^M-\delta p^W,
&
d^{\mathrm B}&\equiv\frac{v'\!\left(P\right)}{\lambda}-\left(1-\Delta\omega^Y\right)F_P,
\label{eq:sp:versions:costates}
\end{align}
with $r\equiv\rho-\dot\lambda/\lambda$ as before. Setting $\Delta=\zeta=0$ and reinstating $\eta$
recovers Section~\ref{sec:sp:problem} exactly.

Five features of this system are worth drawing out.

\smalltitle{(i) The materials arbitrage survives intact.} \eqref{eq:sp:versions:foc:N} has the same
form as~\eqref{eq:sp:solution:arbitrage}: the value of a secondary tonne equals the full cost of a
virgin tonne plus the emission its recovery avoids. All that changes is that extraction resources
are charged at $\left(1+\Delta\omega^N\right)$ rather than at par. The central mechanism of
Section~\ref{sec:sp:solution:recycling} --- scarcity raises $B$, which raises $x$, $\alpha$ and
$\varpi$ --- is unaffected in form, since Proposition~\ref{prop:sp:circularity} used only the
curvature of $a$ and the convexity of $c^T$. The transition is still scarcity-driven; it is now
also, second-order, composition-driven.

\smalltitle{(ii) The wedges are real, not artefacts.} Proposition~\ref{prop:sp:nowedge} said that a
posited reduced form generates spurious taxes on consumption and output. Version A generates wedges
of formally the same shape --- compare~\eqref{eq:sp:versions:foc:C}
with~\eqref{eq:sp:solution:spurious-C} --- but for a different and legitimate reason. In the
reduced form the wedge came from letting the planner reduce economy-wide waste without reducing the
mass drawn in. Here total mass is still conserved; what consumption does is change the \emph{split}
of that mass between the durable stock and the waste flow. That is a real margin, and it should be
priced. The two cases are distinguished by the multiplier: the spurious wedge is proportional to
$p^W$, the genuine one to $\chi\left(p^W-p^M\right)$, and it vanishes when capital is not a sink
($\chi=0$ or $\phi^I=0$) or when postponement is worthless ($p^W=p^M$).

\smalltitle{(iii) The capital wedge is damped.} Comparing~\eqref{eq:sp:versions:foc:I}
with~\eqref{eq:sp:solution:q}, the materials wedge on $q$ is scaled by $\left(1-\chi\right)$. Raising
gross formation embodies more mass in capital, but it also raises $\phi^IG$ in the denominator
of~\eqref{eq:sp:versions:Omega-B}, diluting $\Omega$ and hence $\Phi^I$ itself; $\chi$ is exactly the
fraction of the direct effect that this feedback offsets. The intensity of capital goods
accommodates the planner rather than constraining it, which is the same fact as the disappearance
of~\eqref{eq:sp:problem:mass}.

\smalltitle{(iv) $\lambda$ is no longer marginal utility.} By~\eqref{eq:sp:versions:foc:C},
$u'(C)=\lambda\left(1+\Delta\omega^C\right)$, so the uniform discounting of
Section~\ref{sec:sp:problem:costates} is lost: $r\equiv\rho-\dot\lambda/\lambda$ is still the rate at
which all goods-denominated shadow prices discount, but the consumption Euler equation acquires a
term in the growth of the wedge,
\begin{align}
\sigma\left(C\right)\frac{\dot C}{C}
&= \frac{\left(1-\Delta\omega^Y\right)F_K+\dot q}{q}-\delta-\rho
-\frac{d}{dt}\ln\!\left(1+\Delta\omega^C\right).
\label{eq:sp:versions:euler}
\end{align}
A wedge that is expected to grow --- which is what a rising $p^W$ does --- tilts consumption toward
the present, an effect with no counterpart in Version B.

\smalltitle{(v) The relative intensities now matter.} Every $\omega^j$ appears
in~\eqref{eq:sp:versions:foc:C}--\eqref{eq:sp:versions:costates}, so
Proposition~\ref{prop:sp:irrelevance} fails. This is the substantive difference between the two
closures, and it is a difference about what the calibration exercise \emph{is}. Under Version B the
$\omega^j$ are a measurement bridge: they translate the model's mass flows into published
material-flow accounts and do nothing else. Under Version A they are technology, identified in
principle from the allocation itself. The choice between the closures is therefore not a robustness
check to be run at the end; it determines what the data are being asked to discipline. Neither is
adopted as \emph{the} baseline here.

\begin{remark}[When the two versions coincide]\label{rem:sp:coincide}
$\Delta=0$ iff $\chi\left(p^W-p^M\right)=0$. The two closures therefore deliver identical
allocations in exactly three circumstances: $\phi^I=0$ (capital embodies no material, so $\chi=0$);
$G=0$ (no gross formation, so nothing is embodied); or $p^W=p^M$ (postponing waste is worth nothing,
which by~\eqref{eq:sp:costate:M} requires $r=0$ or $p^W=0$). Otherwise they differ at first order,
and the difference is $O\!\left(\Omega\chi p^W\right)$ --- small when the economy embodies little of
its throughput in durables, large when it embodies much. It is worth reporting $\chi$ alongside
$a\varpi$ for this reason: it is the statistic that says how much the choice of closure can matter.
\end{remark}


\subsection{Simplified waste-generation schemes}\label{sec:sp:versions:simple}

The general scheme of Section~\ref{sec:sp:setup:materialaccounting} lets waste arise from every use
of the final good and, separately, from two physical flows measured in tonnes. Much of it can be
switched off by setting parameters to zero. The switches below are used to build up, test and
decompose the model; the useful question about each is not whether it is realistic but whether it
changes the allocation or only the reported mass flows.

\begin{enumerate}[label=\textbf{(W\arabic*)},ref=W\arabic*,leftmargin=3.2em,itemsep=0.35em]

\item\label{sw:sp:nature} \textbf{No nature-attributable waste:}
$\Omega^{N,S}=\Omega^{D,S}=0$, so $\mathcal N\equiv0$. Overburden and gangue are ignored, and every
tonne of waste has passed through production. The ledger~\eqref{eq:sp:setup:ledger-solved} collapses
to $W=\left(N-\dot M^K\right)/\left(1-a\varpi\right)$; the extraction and exploration
conditions~\eqref{eq:sp:foc:N} and~\eqref{eq:sp:foc:D} lose their $p^W$ terms, giving
$B=C^N_N+p^S+d^Wp^P$ and $p^S+p^X=C^D_D$; and~\eqref{eq:sp:foc:W-solved} becomes explicit rather
than implicit,
\begin{align}
p^W &= c^W\!\left(\varpi\right)+\left(1-\varpi+d^W\varpi\right)p^P
-\varpi\alpha\left(C^N_N+p^S+d^Wp^P\right).
\label{eq:sp:versions:pW-simple}
\end{align}
\emph{Affects the allocation under both versions.}

\item\label{sw:sp:omegaY} \textbf{Waste from output only:}
$\omega^Y=1$ and $\omega^C=\omega^{N}=\omega^{D}=\omega^W=0$, the textbook $W\propto Y$ scheme.

\item\label{sw:sp:omegaYC} \textbf{Waste from output and consumption only:}
$\omega^{N}=\omega^{D}=\omega^W=0$, leaving $\omega^Y,\omega^C$ free.

\item\label{sw:sp:uniform} \textbf{Uniform intensities:} $\omega^j=1$ for every $j$, so that
$\mathcal D$ is total final-good use and $\Omega$ is a single economy-wide intensity.

\item\label{sw:sp:phiI} \textbf{Capital embodies nothing:} $\Phi^I=0$ (Version B) or $\phi^I=0$
(Version A). Then $M^K\equiv0$, the state and its price $p^M$ drop out, $q=1$, the mass
constraint~\eqref{eq:sp:problem:mass} is vacuous, and by Remark~\ref{rem:sp:coincide} the two
versions coincide. The ledger reduces to $W=\left(N+\mathcal N\right)/\left(1-a\varpi\right)$
throughout, not only at $\dot M^K=0$. This is the cleanest available stripped model and the natural
first target for the numerical solver. \emph{Affects the allocation under both versions.}

\item\label{sw:sp:explore} \textbf{No exploration:} $D\equiv0$ and $C^D\equiv0$. The states $X$ and
$S$ reduce to a single depleting reserve, $\dot S=-N$, and $p^X$ drops out; $p^S$ becomes the pure
modified-Hotelling rent~\eqref{eq:sp:costate:S}. \emph{Affects the allocation under both versions},
and note that it does \emph{not} restore a steady state --- see~\ref{sw:sp:CDX}.

\item\label{sw:sp:CDX} \textbf{Exploration costs independent of past discovery:} $C^D_X=0$, so $X$
ceases to be a state and $p^X\equiv0$. The exploration condition becomes
$p^S=C^D_D+\Omega^{D,S}p^W$. This is the switch that repairs
Proposition~\ref{prop:sp:nosteadystate}: with $X$ gone, $\dot S=D-N=0$ is compatible with $D=N>0$,
and a genuine interior steady state exists, characterised by the quasi-stationary
block~\eqref{eq:sp:solution:quasistationary} together with
$N\left(1+\Omega^{N,S}+\Omega^{D,S}\right)=W\left(1-a\varpi\right)$ and $\Phi^K=\Phi^I$.
\emph{Affects the allocation under both versions.}

\item\label{sw:sp:nopoll} \textbf{No pollution channel in production:} $F_P=0$, so damages enter
only through $v(P)$ and $d=v'(P)/u'(C)$. \emph{Affects the allocation under both versions}, though
only through the level of $p^P$.

\end{enumerate}

\smalltitle{The classification.} Switches~\ref{sw:sp:omegaY}--\ref{sw:sp:uniform} are of a different
kind from the rest, and the difference is exactly Proposition~\ref{prop:sp:irrelevance}.

\begin{proposition}[Waste-composition switches are allocationally inert under Version B]
\label{prop:sp:inert}
Under Version B, any restriction on the relative intensities $\left\{\omega^Y,\omega^C,\omega^{N},
\omega^{D},\omega^W\right\}$ --- including~\ref{sw:sp:omegaY}, \ref{sw:sp:omegaYC}
and~\ref{sw:sp:uniform} --- leaves the optimal paths of every state, every control and every shadow
price unchanged. It changes only the derived quantities $\Omega$ and $\Phi^Y$
in~\eqref{eq:sp:setup:Omega}. Under Version A the same restrictions change every condition
in~\eqref{eq:sp:versions:foc:C}--\eqref{eq:sp:versions:costates}.
\end{proposition}

The first part is immediate from Proposition~\ref{prop:sp:irrelevance}; the second from inspection of
$\Delta\omega^j$. Table~\ref{tab:sp:switches} collects the classification.

\begin{table}[H]
\centering
\caption{Simplifications of the waste-generation scheme, and what each affects.}
\label{tab:sp:switches}
\setlength{\tabcolsep}{5pt}
\renewcommand{\arraystretch}{1.05}
\begin{tabular}{@{}p{0.30\textwidth}p{0.20\textwidth}p{0.20\textwidth}p{0.22\textwidth}@{}}
\toprule
Switch & Version B & Version A & States removed \\
\midrule
\ref{sw:sp:nature}\ \ no $\mathcal N$          & allocation & allocation & --- \\
\ref{sw:sp:omegaY}\ \ waste from $Y$ only      & reporting only & allocation & --- \\
\ref{sw:sp:omegaYC}\ \ waste from $Y,C$ only   & reporting only & allocation & --- \\
\ref{sw:sp:uniform}\ \ uniform $\omega^j$      & reporting only & allocation & --- \\
\ref{sw:sp:phiI}\ \ $\Phi^I=0$                 & allocation & allocation (collapses to A) & $M^K$ \\
\ref{sw:sp:explore}\ \ no exploration          & allocation & allocation & $X$ \\
\ref{sw:sp:CDX}\ \ $C^D_X=0$                   & allocation & allocation & $X$ \\
\ref{sw:sp:nopoll}\ \ $F_P=0$                  & allocation & allocation & --- \\
\bottomrule
\end{tabular}
\end{table}

Two readings of the table matter for what follows. First, the only simplifications that are free
under Version B are precisely the ones about the \emph{composition} of waste across final-good uses
--- the part of the accounting that is hardest to discipline empirically and that occupies most of
Section~\ref{sec:sp:setup:materialaccounting}. Under Version B that section can be read as
bookkeeping; under Version A it cannot. Second, none of the switches removes the ledger itself.
Mass conservation is not a modelling option in this framework, and no setting of the parameters
recovers a specification in which recycling avoids waste rather than deferring it.
